\section{Proofs}\label{app:proofs}

Notation: $m = \nmin$; $\mathcal N(\mu,1)$ is the normal law; $\varphi$ and $\Phi$ are the standard normal density and distribution function; $\psi(u) = \varphi(\Phi^{-1}(u))$; $\|\cdot\|$ is the supremum norm.

\subsection{Proof of \texorpdfstring{\cref{prop:toll}}{the intercept toll}}\label{app:proof-toll}

For a child $c$ let $D_c = G_c - rH_c = -(k_c - \phat\,n_c)$. We first show the exact identity
\begin{equation}\label{eq:toll}
  \Gain = \sum_{c\in\{L,R\}} \frac{(D_c - r\lambda)^2}{H_c + \lambda} - \lambda r^2\,\frac{H + 2\lambda}{H+\lambda}.
\end{equation}
Write $G_c = rH_c + D_c = r(H_c+\lambda) + (D_c - r\lambda)$ and expand: $G_c^2/(H_c+\lambda) = r^2(H_c+\lambda) + 2r(D_c - r\lambda) + (D_c - r\lambda)^2/(H_c+\lambda)$. Summing over the two children and using $D_L + D_R = 0$ and $H_L + H_R = H$ gives $r^2H - 2\lambda r^2 + \sum_c (D_c-r\lambda)^2/(H_c+\lambda)$. The parent term is $(rH)^2/(H+\lambda) = r^2H - \lambda r^2H/(H+\lambda)$. Subtracting gives \eqref{eq:toll}.

For the approximation, $\sum_c D_c^2/H_c$ is the right-hand side of \cref{prop:chi2} at the root, $\kappa X^2$ with $\kappa = \phat(1-\phat)/(b(1-b))$. Expanding the squares in \eqref{eq:toll}, the gain minus $(\kappa X^2 - \lambda r^2)$ is the sum of four terms: $-\sum_c D_c^2\lambda/(H_c(H_c+\lambda))$, $-2r\lambda\sum_c D_c/(H_c+\lambda)$, $r^2\lambda^2\sum_c 1/(H_c+\lambda)$ and $-\lambda^2r^2/(H+\lambda)$. With $H_{\min} = \min(H_L, H_R)$, and since $D_c/H_c = (\phat - \phat_c)/(b(1-b))$ and $r = (b-\phat)/(b(1-b))$, bounding each term gives
\begin{equation}\label{eq:remainder}
  \bigl|\Gain - (\kappa X^2 - \lambda r^2)\bigr| \le \frac{\lambda}{H_{\min}}\bigl(\kappa X^2 + 2\lambda r^2\bigr) + 2\lambda r^2\sum_c \frac{|\phat_c - \phat|}{|b-\phat|},
\end{equation}
which is small when $H_c \gg \lambda$ and $|\phat_c - \phat| \ll |b - \phat|$. Solving $\kappa X^2 - \lambda r^2 > 0$ for $X^2$ gives \eqref{eq:threshold}; the asymptotic forms use $b \to \sigma(-2)$, $b/(1-b) \to e^{-2}$ (version 2.x) and $b = 1/2$ (version 1.x). \hfill$\square$

\subsection{Proof of \texorpdfstring{\cref{thm:limit}}{Theorem 1}}\label{app:proof-limit}

\emph{Coupling.} The claim is about laws, so we may build all sample sizes on one probability space. For one feature, fix the $m$ positive values $Z_1,\dots,Z_m$, i.i.d.\ with distribution $F_1$, and an infinite i.i.d.\ sequence $\xi_1, \xi_2, \dots$ with distribution $F_0$; the sample of size $n$ consists of the positives and $\xi_1, \dots, \xi_{n-m}$. Different features use independent copies, which gives independence across features at every $n$ because features are conditionally independent given $Y$. We show almost-sure convergence, which implies convergence in distribution.

\emph{Rank fractions.} Let $F^-_n$ be the empirical distribution function of $\xi_1,\dots,\xi_{n-m}$. The rank fraction of positive $i$ among all $n$ values is $\rho_i^n = n^{-1}[(n-m)F^-_n(Z_i) + \#\{l \le m : Z_l \le Z_i\}]$, so
\[
  |\rho^n_i - F_0(Z_i)| \le \|F^-_n - F_0\| + 2m/n \to 0 \quad\text{a.s.}
\]
by the Glivenko--Cantelli theorem. Let $V_i = F_0(Z_i)$. Because $V$ has no atoms, almost surely the $V_i$ are distinct, lie in $(0,1)$ and differ from $a = w/m$ and $1-a$.

\emph{The counting process.} Exact greedy search evaluates, for $j = 1,\dots,n-1$, the split that sends the $j$ smallest values left, so $u = j/n$ and $k_L(j/n) = N_n(j/n)$ with $N_n(u) = \#\{i : \rho_i^n \le u\}$. Fix $\omega$ outside a null set and $n$ large enough that the $\rho^n_i$ are in the same order as the $V_i$. Let $\ell_n$ be the increasing piecewise-linear bijection of $[0,1]$ with $\ell_n(0) = 0$, $\ell_n(1) = 1$ and $\ell_n(\rho^n_i) = V_i$. Then $N_n = N\circ\ell_n$ exactly, with $N(u) = \#\{i : V_i \le u\}$, and $\|\ell_n - \mathrm{id}\| \le \max_i|\rho^n_i - V_i| \to 0$.

\emph{The bracket.} With the prevalence intercept, $H_L = n_L\phat(1-\phat) = mu(1-\phat)$, so the admissible thresholds are $u \in [a_n, 1-a_n]$ with $a_n = w/(m(1-\phat)) \to a$, and the bracket of \eqref{eq:process} is $\beta_n(u) = [mu(1-\phat)+\lambda]^{-1} + [m(1-u)(1-\phat)+\lambda]^{-1}$. Since $w > 0$, $\beta_n \to \beta$ uniformly on $[a/2, 1-a/2]$, where $\beta$ is the bracket of \eqref{eq:limit}, and $\beta$ is uniformly continuous there (this covers $\lambda = 0$).

\emph{Process convergence.} Extend $\Gamma$ by \eqref{eq:limit} to $[a/2, 1-a/2]$. Since $N_n = N\circ\ell_n$ with $\|\ell_n - \mathrm{id}\|\to 0$, the counting processes converge in $D[0,1]$ with the $J_1$ topology \citep{billingsley1999convergence}. For the gains, on $[a_n, 1-a_n]$, $\Gain_n(u) - \Gamma(\ell_n(u)) = (N(\ell_n u) - mu)^2\beta_n(u) - (N(\ell_n u) - m\ell_n u)^2\beta(\ell_n u)$. Since $0 \le N \le m$, $|mu - m\ell_n u| \le m\|\ell_n - \mathrm{id}\|$, $\beta_n \to \beta$ uniformly and $\beta$ is uniformly continuous, $\sup_{u\in[a_n,1-a_n]}|\Gain_n(u) - \Gamma(\ell_n(u))| \to 0$.

\emph{Suprema.} The search maximises over the grid $T_n = \{j/n\}\cap[a_n, 1-a_n]$, whose image $\ell_n(T_n)$ has mesh tending to zero and end points tending to $a$ and $1-a$. $\Gamma$ is continuous except for finitely many jumps at the $V_i$, none at $a$ or $1-a$. Lower bound: the supremum of $\Gamma$ over $[a,1-a]$ is a value at a continuity point or a one-sided limit at a jump, and every open interval inside $[a,1-a]$ eventually contains points of $\ell_n(T_n)$, so $\liminf_n\max_{T_n}\Gain_n \ge \sup_{[a,1-a]}\Gamma$. Upper bound: for any $\eta > 0$, eventually $\ell_n(T_n)\subset[a-\eta, 1-a+\eta]$, so $\limsup_n\max_{T_n}\Gain_n \le \sup_{[a-\eta,1-a+\eta]}\Gamma$; let $\eta\to 0$ using the continuity of $\Gamma$ at $a$ and $1-a$. Hence $\max_{T_n}\Gain_n \to \sup_{[a,1-a]}\Gamma$ almost surely, jointly over the independently generated features, and the law of $\Gamma$ is a function of $(m, G, \lambda, w)$ only. \hfill$\square$

\begin{lemma}[shadows in the limit]\label{lem:shadow}
Let a shadow be a uniformly random row-permutation, independent of the data, of a candidate column whose negatives are i.i.d.\ from a continuous $F_0$. In the coupling above, the shadow's placements converge almost surely to i.i.d.\ uniform variables independent of the limits of all candidates and of other shadows, so its gain process converges to that of a null feature, independently of the others.
\end{lemma}

\begin{proof}
With probability at least $1 - m^2/(n-m)$ the permutation sends every positive row to a row that is negative in the source column, so the shadow's $m$ ``positive'' values are $m$ draws without replacement from the source's i.i.d.\ $F_0$ negatives, at positions chosen independently of the data; they are therefore i.i.d.\ $F_0$ and independent of the source's positives. The shadow's ``negative'' values are the source's values with $m$ of the negatives replaced by its $m$ positives, so their empirical distribution function is within $2m/(n-m)$ of the source's negative one in the supremum norm and converges uniformly to $F_0$. The argument of the proof above then gives placements $F_0(\xi)$, i.i.d.\ uniform. Distinct shadows use independent permutations, so their selected rows are disjoint with probability tending to one, and their limits are independent of each other and of the candidates' positives.
\end{proof}

\subsection{Proof of \texorpdfstring{\cref{prop:fixed}}{Proposition 3}}\label{app:proof-fixed}

$N(u)$ counts the i.i.d.\ placements with $V_i \le u$, each with probability $G(u)$, so it is binomial. $\E(N - mu)^2 = mG(1-G) + m^2(G-u)^2$; dividing by $mu(1-u)$ gives the mean. If $\sqrt m(G_m(u) - u) \to -c$ then $G_m(u)\to u \in (0,1)$, so $(N - mG_m)/\sqrt{mG_m(1-G_m)} \Rightarrow \mathcal N(0,1)$ by the de Moivre--Laplace theorem for triangular arrays, and $(mG_m - mu)/\sqrt{mu(1-u)} \to -c/\sqrt{u(1-u)}$.

For \eqref{eq:twopi}, write $x = \Phi^{-1}(u)$ and $f_\delta(x) = (\int_0^1\varphi(x-\theta\delta)\,d\theta)^2/(\Phi(x)(1-\Phi(x)))$, so that $\nu(u)/(m\delta^2) = f_\delta(x)$ by the mean value theorem. For $|\delta|\le 1$, $\varphi(x-\theta\delta) \le \varphi(x)e^{|x|}$, so $f_\delta \le g$ with $g(x) = \varphi(x)^2e^{2|x|}/(\Phi(x)(1-\Phi(x)))$, and by symmetry and the Mills-ratio bound $1-\Phi(x) \ge \varphi(x)x/(1+x^2)$, $g(x)\to 0$ as $|x|\to\infty$. On compacts $f_\delta \to f_0$ uniformly, so $\sup f_\delta \to \sup f_0$ with $f_0(x) = \varphi(x)^2/(\Phi(x)(1-\Phi(x)))$. For independent standard normals $Z_1, Z_2$ and $x\neq 0$, $(2\Phi(x)-1)^2 = \P(|Z_1|\le|x|, |Z_2|\le|x|) < \P(Z_1^2+Z_2^2 \le 2x^2) = 1 - e^{-x^2}$, because the square lies strictly inside the disc of radius $|x|\sqrt2$. Hence $\Phi(x)(1-\Phi(x)) = \frac14[1-(2\Phi(x)-1)^2] > \frac14 e^{-x^2}$ and $f_0(x) < (e^{-x^2}/2\pi)/(e^{-x^2}/4) = 2/\pi = f_0(0)$. As $f_0$ is continuous with a unique maximiser and $f_\delta\to f_0$ uniformly on compacts under the envelope $g$, every maximiser of $f_\delta$ tends to $0$, that is $u\to 1/2$. \hfill$\square$

\subsection{Proof of \texorpdfstring{\cref{thm:scaling}}{Theorem 2}}\label{app:proof-scaling}

Fix one feature with shift $\delta = \delta_m$ and $s_m = m\delta_m^2\to s$. Let $U_i = \Phi(Z_i - \delta)$, i.i.d.\ uniform. Then $\Phi(Z_i)\le u$ if and only if $U_i \le q_\delta(u) = \Phi(\Phi^{-1}(u)-\delta)$, so $N(u) = m\hat G_m(q_\delta(u))$ with $\hat G_m$ the empirical distribution function of the $U_i$. With the uniform empirical process $\alpha_m(t) = \sqrt m(\hat G_m(t) - t)$,
\[
  \frac{N(u) - mu}{\sqrt m} = \alpha_m(q_\delta(u)) + \sqrt m\,(q_\delta(u) - u).
\]
(a) \emph{Drift.} With $x = \Phi^{-1}(u)$, $u - q_\delta(u) = \delta\int_0^1\varphi(x-\theta\delta)\,d\theta$ and $|\varphi'|\le\varphi(1)$, so $|\sqrt m(u-q_\delta(u)) - \sqrt m\delta\psi(u)| \le \varphi(1)s_m/\sqrt m \to 0$ uniformly in $u$, and the drift converges uniformly to $-\sqrt s\,\psi(u)$.
(b) \emph{Noise.} $\|q_\delta - \mathrm{id}\| \le \varphi(0)\delta\to 0$. The uniform empirical process converges to a Brownian bridge $\mathbb B$ in $D[0,1]$ and is asymptotically equicontinuous \citep{shorack1986empirical,billingsley1999convergence}, so $\|\alpha_m\circ q_\delta - \alpha_m\|\to 0$ in probability and $\alpha_m\circ q_\delta\Rightarrow\mathbb B$; since the limit is continuous, the convergence holds in the supremum norm.
(c) \emph{Bracket.} On $[\varepsilon,1-\varepsilon]$, $m[(mu+\lambda)^{-1} + (m(1-u)+\lambda)^{-1}] \to 1/(u(1-u))$ uniformly. Combining (a)--(c) with the continuous mapping theorem, $\Gamma_{\delta_m}\Rightarrow\Gamma^{(s)}$ uniformly on the window; independent features give independent limits.
(d) \emph{Win probability.} $f\mapsto\sup_{[\varepsilon,1-\varepsilon]}f$ is continuous for the supremum norm, so the suprema converge jointly. For $s = 0$, $\sup\Gamma^{(0)} = (\sup_u|\mathbb B(u)|/\sqrt{u(1-u)})^2$. By Tsirelson's theorem \citep{tsirelson1975density,lifshits1995gaussian} applied to the symmetrised centred Gaussian family $\{\pm\mathbb B(u)/\sqrt{u(1-u)}\}$, the law of its supremum has no atoms except possibly at the left end of its support, which is $0$ and carries no mass. So the null suprema have continuous laws, ties have probability zero, and the win probability converges by the portmanteau theorem.
(e) \emph{Monotonicity.} For $c > 0$ let $A_c = \{f\in C[\varepsilon,1-\varepsilon] : |f(u)|\le\sqrt{c\,u(1-u)}\ \forall u\}$, a closed, convex, symmetric set, so $\P(\sup\Gamma^{(s)}\le c) = \P(\mathbb B - \sqrt s\psi\in A_c)$. The law $\gamma$ of $\mathbb B$ is a centred Gaussian Radon measure, hence log-concave \citep{borell1974convex,bogachev1998gaussian}, and $t\mapsto\gamma(A_c + t\psi)$ is log-concave and even, hence non-increasing for $t\ge 0$ (Anderson's inequality, \citealp{anderson1955integral}). So $\sup\Gamma^{(s)}$ is stochastically non-decreasing in $s$, and $\omega_\varepsilon(s,K)$ is non-decreasing in $s$.
(f) For $s = 0$ the $K+1$ suprema are i.i.d.\ with a continuous law, so each is the largest with probability $1/(K+1)$.

\emph{General local alternatives.} The argument uses the Gaussian model only in (a). If instead $\sqrt m(u - G_m(u))\to h(u)$ uniformly on the window for a continuous $h$, represent the placements as $V_i = G_m^{-1}(U_i)$ with $U_i$ i.i.d.\ uniform, so that $V_i\le u$ if and only if $U_i\le G_m(u)$ and $N(u) = m\hat G_m(G_m(u))$, with $\hat G_m$ the empirical distribution function of the $U_i$. Then $(N(u)-mu)/\sqrt m = \alpha_m(G_m(u)) + \sqrt m(G_m(u)-u)$; step (b) needs only $\|G_m - \mathrm{id}\|\to 0$, which holds on the window, and the limit is $(\mathbb B(u) - h(u))^2/(u(1-u))$. Steps (d)--(f) hold with $\sqrt s\,\psi$ replaced by $h$, and (e) gives monotonicity along $h \mapsto t h$, $t \ge 0$. \hfill$\square$

\subsection{Proof of \texorpdfstring{\cref{prop:level}}{Proposition 4}}\label{app:proof-level}

By (A0), (A3) and the strong null, the rank vector of $X_j$ is uniform over permutations and independent of $Y$ and the other candidates; by (A1) the importances depend on the columns only through such rank vectors. (A3) is needed only for $X_j$ and the shadows' source columns. A shadow's rank vector is the rank vector of its source composed with an independent uniform permutation, so it is again uniform and independent of everything else, as is the rank vector of an i.i.d.\ continuous noise column. So the rank vectors of $X_j$ and of the $K$ shadows are i.i.d.\ and independent of $Y$ and the remaining columns, and by (A2) the $K+1$ importances are exchangeable; each is the strict maximum with probability at most $1/(K+1)$. For per-feature shadows, (A0) and the strong null give $X_j = \tau x$ for its multiset of values $x$ and a uniform permutation $\tau$ independent of $(Y, X_{C\setminus j})$; the shadows are $\tau_k x$ with $\tau_k = \varrho_k\tau$ for the shadows' independent uniform permutations $\varrho_k$, and $\tau, \tau_1, \dots, \tau_K$ are i.i.d.\ uniform. The $K+1$ columns are therefore exchangeable given the rest, and (A2) gives exchangeable importances without (A1) or (A3); ties, as with discrete columns, only make the inequality strict. Shadows of other candidates do no harm, being independent of the $\varrho_k$. \hfill$\square$

\subsection{Proof of \texorpdfstring{\cref{prop:rounds}}{Proposition 5}}\label{app:proof-rounds}

Given $D$ the rounds are i.i.d., which gives the binomial law, and $\E\,\theta_j(D) = \P(\text{hit}) \le 1/(K+1)$ by \cref{prop:level} applied to each round (fresh shadows and a fresh column order make each round satisfy its conditions). The bound in \eqref{eq:rounds} is Markov's inequality for $T_j$: $\P(T_j \ge r_R+1)\le \E T_j/(r_R+1)$. As $R\to\infty$, $r_R/R\to 1/2$ and $T_j/R\to\theta_j(D)$ almost surely given $D$; for $\theta < 1/2$, $\P(\Bin(R,\theta) > r_R)\to 0$, so by dominated convergence $\limsup_R\P(\text{keep})\le\P(\theta_j(D)\ge 1/2)$. Markov's inequality for $\theta_j^2$ gives $\P(\theta_j\ge 1/2)\le 4\E\theta_j^2 = 4\operatorname{Var}\theta_j + 4(\E\theta_j)^2$, and Markov's inequality for $\theta_j$ gives $\P(\theta_j\ge 1/2)\le 2\E\theta_j\le 2/(K+1)$. \hfill$\square$

\subsection{Proof of \texorpdfstring{\cref{prop:mw}}{Proposition 6} and \texorpdfstring{\cref{cor:efficiency}}{Corollary 2}}\label{app:proof-mw}

Without ties, $U$ counts the (positive, negative) pairs with the positive larger, so $W_n = m^{-1}\sum_i F^-_n(Z_i)$ with $F^-_n$ as in \cref{app:proof-limit}. By Glivenko--Cantelli $F^-_n\to F_0$ uniformly, which gives the first form of \eqref{eq:mw}; the second follows from $\int_0^1 \mathbf 1\{V_i\le u\}\,du = 1 - V_i$.

For the normal limit, the $V_i$ form a triangular array of i.i.d.\ variables in $[0,1]$ with mean $\AUC_m$. Since the placement laws $G_m\to\mathrm{id}$ uniformly, $V_i$ converges in law to a uniform variable, and by bounded convergence $\operatorname{Var}V_i \to 1/12$. The Lindeberg--Feller theorem (the summands are bounded and their variance is bounded below) gives $\sqrt{12m}\,(W - \AUC_m)\Rightarrow \mathcal N(0,1)$. Finally $\sqrt{12m}\,(\AUC_m - \tfrac12) = \sqrt{12\,\sauc/(4\pi)}\to\sqrt{3s_*/\pi}$, and the power of the two-sided test with the null standardisation follows.

For \cref{cor:efficiency}, let $X^- \sim f$ and $X^+ = X' + \delta$ with $X'\sim f$ independent. At a fixed $u$ with $x = F^{-1}(u)$ a continuity point of $f$, $u - G_\delta(u) = F(x) - F(x-\delta) = \delta J_f(u) + o(\delta)$, so by \cref{prop:fixed} the squared mean of the standardised count is $s J_f(u)^2/(u(1-u))$ in the limit. For the rank test, $\AUC - \tfrac12 = \int [F(x+\delta) - F(x)] f(x)\,dx = \delta\int f^2 + o(\delta)$ by dominated convergence (as $f$ is bounded), so $\sauc = 4\pi s(\int f^2)^2(1+o(1))$ and its squared mean is $3\sauc/\pi = 12 s(\int f^2)^2$. Both statistics are asymptotically normal with unit variance, so the ratio of squared means is the Pitman efficiency. The numerical values in \eqref{eq:efficiency} are computed on a grid of $2\times 10^5$ points in $u$. For the normal, $J = \psi$, $\sup_u\psi^2/(u(1-u)) = 2/\pi$ and $\int\varphi^2 = 1/(2\sqrt\pi)$, so $\mathrm{ARE} = (2/\pi)/(12/(4\pi)) = 2/3$. \hfill$\square$
